\subsection{From logic-based agents to BDI}
\label{sec:logic-based-agents}
\slidesrc{BDI architecture: 1--6}

The deck opens with the two definitions of agent already seen in
Section~\ref{sec:what-is-an-agent}:
\begin{itemize}[nosep]
  \item \textbf{weak notion} (Jennings, Sycara, Wooldridge 1998): situated, autonomous,
        flexible (reactive, proactive, social);
  \item \textbf{strong notion}: a computer system either \emph{conceptualised} or
        \emph{implemented} using concepts that are more usually applied to humans:
        \emph{mentalistic notions} (Shoham, \emph{Agent-oriented programming}, 1993) and
        \emph{emotional notions} (Bates, \emph{The role of emotion in believable agents},
        1994).
\end{itemize}

\paragraph{Logic-based agents (slide 5).} The simplest way to give an agent ``a mind'' is to
give it a logic program: knowledge as clauses, reasoning as proof.
\begin{minted}[linenos=false]{prolog}
mortal(X) :- human(X).
human(socrates).
?- mortal(socrates).
yes
\end{minted}
This is the ``traditional'' approach to building intelligent systems (EMAS keynote, slide
32): \textbf{logical formulae} give a symbolic representation of the agent's environment and
desired behaviour, and \textbf{logical deduction} (theorem proving) manipulates it. A pure
theorem prover, however, only answers questions; an agent must also \emph{act}, pursue goals
over time and react to changes. This is what the BDI architecture adds.

\begin{nbfact}{AgentSpeak(L) (slide 6)}
The \textbf{BDI architecture} is one of the best known and most studied architectures for
\textbf{cognitive agents}. \textbf{AgentSpeak(L)} is an elegant, logic-based programming
language inspired by the BDI architecture (A.~S.~Rao, ``AgentSpeak(L): BDI agents speak out
in a logical computable language'', MAAMAW 1996).
\end{nbfact}
